\chapter{Emulated platforms and chips}

\begin{longtable}{@{}p{0.22\textwidth} p{0.27\textwidth} p{0.43\textwidth}@{}}
\toprule \textbf{Platform} & \textbf{Chip} & \textbf{Synthesis} \\ \midrule \endhead
ZX Spectrum 128, Amstrad CPC, Atari ST, MSX & AY-3-8910 / YM2149 & 3 square voices, noise, hardware envelope \\
ZX Spectrum (TurboSound) & 2 $\times$ AY & 6 voices \\
ZX Spectrum (TurboFM) & 2 $\times$ YM2203 & 4-operator FM + AY \\
SAM Coupé & SAA1099 & 6 square voices, noise, envelopes \\
Commodore 64 & MOS 6581 / 8580 (SID), up to 3 chips & 3 oscillators, analogue filter, ring modulation, sync \\
NES / Famicom & Ricoh 2A03 / 2A07 & 2 pulses, triangle, noise, DPCM \\
Super Nintendo & SPC700 + S-DSP & 8 BRR sample voices, ADSR, FIR echo \\
Game Boy & LR35902 (APU) & 2 pulses, programmable wave, noise \\
PC Engine / TurboGrafx-16 & HuC6280 & 6 wavetable voices 32 $\times$ 5 bits, noise, DDA, LFO \\
Atari 8-bit & POKEY & 4 voices, noise polynomials, filters \\
Mega Drive / Genesis & YM2612 + SN76489 & 6-voice 4-operator FM, DAC, PSG \\
PC (AdLib, Sound Blaster) & YM3526 / YM3812 / YMF262 & 2- or 4-operator FM, rhythm mode \\
Sharp X68000, arcade & YM2151 & 8-voice 4-operator FM \\
Sega Model 1 and 2 arcade & YMW258 (MultiPCM) & 28 sample voices \\
Sega Saturn, Model 2A/3 arcade & YMF292 (SCSP) & 32 PCM/FM voices, effects DSP \\
Roland MT-32 & LA (Munt) & linear arithmetic synthesis \\
\bottomrule
\end{longtable}

\section{AY-3-8910 and YM2149 (ZX Spectrum, Amstrad CPC, Atari ST, MSX)}
The AY-3-8910 programmable sound generator was designed by General Instrument in the late 1970s. Yamaha produced a
compatible variant, the YM2149, used in the Atari ST. It is found in the ZX Spectrum 128, Amstrad CPC, MSX, Oric,
many arcade boards and expansion cards.

Voice frequency: $f = F_{clock} / (16 \times N)$, where $N$ is the 12-bit period. Volume is 4 bits (logarithmic, about
3~dB per step on the AY, finer steps on the YM2149). The envelope has a 16-bit period and eight shapes. The clock
depends on the machine: 1.7734~MHz for the ZX Spectrum, 1.75~MHz for the Pentagon, 2~MHz for the Atari ST, 1~MHz for
the Amstrad CPC. Composers built chords from fast arpeggios and used the envelope at audible rates as a fourth
waveform.

\soft reads the formats of the scene trackers (Pro Tracker, Sound Tracker, ASC, SQ-Tracker, Vortex Tracker\dots),
register logs (PSG, YM, VTX) and AY files containing the original player code, run on an emulated Z80.

\section{TurboSound and TurboFM}
Russian ZX Spectrum clones received expansion cards: TurboSound doubles the AY (six voices), TurboFM adds two Yamaha
YM2203 chips, each combining three FM voices and an AY. \soft then shows two columns of voices and the FM operators.

\section{SAA1099 (SAM Coupé)}
The Philips SAA1099 is used in the SAM Coupé, the British successor of the Spectrum, and in Creative Music System
cards. It produces six square channels, two noise generators and two envelopes, with per-channel stereo. E-Tracker is
the main format.

\section{SID (Commodore 64)}
The SID (Sound Interface Device) was designed by Bob Yannes at MOS Technology for the Commodore~64, released in 1982.
Designed as a real synthesizer, it was far ahead of other chips of its time and gave rise to a generation of
composers: Rob Hubbard, Martin Galway, Ben Daglish, Jeroen Tel, Chris Hülsbeck and many others.

Voice frequency: $f = N \times F_{clock} / 2^{24}$, with a 16-bit $N$; the clock is about 0.985~MHz in PAL and
1.023~MHz in NTSC. The pulse width is 12 bits. Sync resets the oscillator at the rate of the previous voice; ring
modulation multiplies the triangle by the previous voice's oscillator (V1 from V3, V2 from V1, V3 from V2). Register
\$D418 sets the 4-bit master volume and the filter mode.

\subsection{6581 and 8580}
The original 6581 has a filter that varies a lot from chip to chip and a current leak producing a click on each volume
change. The 8580, its replacement from 1986, has a more regular filter but sounds different; it strongly attenuates
samples played through the volume. \soft emulates both and offers filter curves measured on real chips.

\subsection{Samples on the SID}
The SID has no sample player, but composers invented two techniques:
\begin{description}
  \item[Volume register] writing the master volume (\$D418) thousands of times per second produces a 4-bit sample;
        it adds to the output, outside the three voices.
  \item[Test bit] freezing and restarting an oscillator lets a value be written into it; the sample then goes through
        a voice and its filter.
\end{description}
The samples lane of the oscilloscope shows the first technique separately.

\subsection{Several SIDs}
Some recent tunes use two or three chips (expansion cartridges, emulators). \soft places them left, right and centre,
or mixes them in mono depending on the setting.

\section{NES and Famicom (Ricoh 2A03)}
The NES processor, the Ricoh 2A03 (2A07 in Europe), contains its sound unit: two pulse channels with selectable duty
cycle (12.5~\%, 25~\%, 50~\%, 75~\%), a triangle, a noise channel and a DPCM channel playing 1-bit samples. Koji Kondo,
Hirokazu Tanaka and Tim Follin wrote classic soundtracks for it. \soft reads NSF and NSFe files, which contain the
game's player code.

\section{Super Nintendo (SPC700 and S-DSP)}
Super Nintendo sound is handled by a subsystem designed by Sony: an SPC700 processor and a signal processor, the
S-DSP, sharing 64~KB of memory. The S-DSP plays eight voices of BRR-compressed samples, with ADSR or GAIN envelope,
stereo volume and a FIR-filtered echo that gives the console's music its depth. It does not synthesize: it plays back
short instrument recordings, so sound quality depends on the samples chosen and on very limited memory.

An SPC file is a snapshot of the sound subsystem memory and registers; an SNSF file also holds the main processor
program, for games driving sound from it.

\section{Game Boy}
The Game Boy has two pulse channels (one with frequency sweep), a programmable wave channel of 32 4-bit steps and a
noise channel. GBS files contain the player code.

\section{PC Engine (HuC6280)}
NEC and Hudson Soft's PC Engine, called TurboGrafx-16 in the United States, integrates sound into its HuC6280
processor. Each of its six channels plays a 32-step 5-bit wavetable written by the program: this is wavetable
synthesis, where the waveform itself is the instrument. Channels 5 and 6 can produce noise; all can switch to DDA
mode, where the processor writes samples directly; channel 2 can modulate the frequency of channel 1 (LFO).

Registers are at \$0800: \$800 selects the channel, \$801 sets left and right master volume, \$802 and \$803 the
12-bit period, \$804 enable, DDA mode and 5-bit volume, \$805 left and right volume, \$806 wave data, \$807 noise,
\$808 and \$809 the LFO. Voice frequency: $f = 3\,579\,545 / (32 \times N)$.

\section{Atari 8-bit (POKEY)}
POKEY, designed by Doug Neubauer at Atari, is used in the Atari 400, 800, XL and XE and in many arcade machines. Its
four channels combine frequency dividers and noise polynomials, giving a very distinctive timbre; two channels can be
paired for better pitch precision. SAP files are played by the ASAP library.

\section{Mega Drive (YM2612 and SN76489)}
Sega's Mega Drive pairs a Yamaha YM2612 FM chip (six 4-operator voices, one able to play 8-bit samples) with the
SN76489 generator inherited from the Master System (three square voices and noise). Yuzo Koshiro, Masato Nakamura and
Motoi Sakuraba used all of its resources. Music is played from VGM logs.

\section{AdLib and Sound Blaster (OPL, OPL2, OPL3)}
On PC, AdLib then Sound Blaster cards use the Yamaha YM3812 (OPL2) then YMF262 (OPL3) FM chips. The OPL2 offers nine
2-operator voices, or six voices and five percussion sounds in rhythm mode; the OPL3 doubles the voice count, adds
4-operator voices and stereo. Since rhythm mode is rarely used, \soft shows the nine melodic voices in a 3~$\times$~3
grid by default.

\section{YM2151 (OPM)}
The OPM offers eight 4-operator voices and a rich LFO. It is found in the Sharp X68000, Yamaha synthesizers of the same
generation and many arcade boards.

\section{Sega Model 1, 2, 2A and 3 arcade boards}
Sega's 3D arcade boards use sample-based chips:
\begin{description}
  \item[Model 1 and Model 2] 68000 processor, YM3438 and two YMW258 (MultiPCM) with 28 sample voices each;
  \item[Model 2A, 2B, 2C and Model 3] SCSP (YMF292), the Saturn chip: 32 PCM voices that can FM-modulate each other,
        envelopes, LFO and a programmable effects DSP for reverb.
\end{description}

\section{Roland MT-32}
Roland's MT-32 MIDI module, widely used by late-1980s PC games, combines samples and synthesis in its LA (linear
arithmetic) synthesis. \soft emulates it with the Munt library, which requires images of the module's ROMs.
